\documentclass[12pt]{article}
\input{litweb-latex/litweb.sty}
\title{LaTeX \# \& \_ Demonstration}
\author{}
\date{}
\LitwebDeclareMiniIndexMeaning{1}{\LitwebMiniIndexName{Graph}: struct, \hyperlink{litweb-1-2}{2}.}
\LitwebDeclareMiniIndexMeaning{2}{\LitwebMiniIndexName{graph}: parameter, \hyperlink{litweb-1-3}{3}.}
\LitwebDeclareMiniIndexMeaning{3}{\LitwebMiniIndexName{graph\_value}: function, \hyperlink{litweb-1-3}{3}.}
\LitwebDeclareMiniIndexMeaning{4}{\LitwebMiniIndexName{main}: function, \hyperlink{litweb-1-1}{1}.}
\LitwebDeclareMiniIndexMeaning{5}{\LitwebMiniIndexName{path}: local variable, \hyperlink{litweb-1-1}{1}.}
\LitwebDeclareMiniIndexMeaning{6}{\LitwebMiniIndexName{squared\_distance}: function, \hyperlink{litweb-1-2}{2}.}
\LitwebDeclareMiniIndexMeaning{7}{\LitwebMiniIndexName{value\_1}: field, \hyperlink{litweb-1-2}{2}.}
\LitwebDeclareMiniIndexMeaning{8}{\LitwebMiniIndexName{x}: parameter, \hyperlink{litweb-1-2}{2}.}
\LitwebDeclareMiniIndexMeaning{9}{\LitwebMiniIndexName{y}: parameter, \hyperlink{litweb-1-2}{2}.}
\begin{document}
\maketitle
\LitwebMiniIndexBegin
\LitwebSection{litweb-1-1}{1}\LitwebSectionTitle{Distance \& escaping.}
\par This \textbf{small \href{\detokenize{https://example.com/a_b?x=1&y=2}}{guide}
matters}, as does \emph{multiline
emphasis}. Literal code preserves \texttt{\ \ @\{not\ active\}\ \ }, while
\texttt{\LitwebMiniUse{7}value\_1\ \&\ value\_2} is Rust. The equation \(d^2 = x^2 + y^2\) remains TeX.
The implementation is in \(\langle\)\texttt{Compute\_value\ \#1}\ \hyperlink{litweb-1-2}{2}\(\rangle\). An [unsafe](javascript:alert)
target stays visible. Unicode prose includes café and an em dash — safely.

\begin{itemize}
\item A list item uses \texttt{\LitwebMiniUse{1}Graph}.
\item A second item continues
on another physical line.
\end{itemize}
\begin{enumerate}
\setcounter{enumi}{2}
\item Ordered values may begin elsewhere.
\item Their numbering continues.
\end{enumerate}
\begin{LitwebCode}
literal \LitwebBackslash{}command \LitwebLeftBrace{} braces \LitwebRightBrace{} % # $ & _ ^ ~
This deliberately long fixed-width example line contains enough words and symbols to exceed an ordinary printed measure without disappearing beyond the right edge.
\end{LitwebCode}
\begin{center}
\LitwebMiniBlockUse{1}\nobreak
\begin{tabularx}{\linewidth}{LCR}
\toprule
Form & Meaning & Count \\
\midrule
\texttt{a\_b} & \emph{centered} & 2 \\
\texttt{Graph} & \href{\detokenize{guide.html}}{safe} & \(n+1\) \\
\bottomrule
\end{tabularx}
\end{center}
\[
\begin{aligned}
d^2 &= x^2 + y^2 \\
    &= 25
\end{aligned}
\]
\LitwebCodeTitle{\(\langle\)\textbf{\texttt{demo.rs}}\ \hyperlink{litweb-1-1}{1}\(\rangle\) \(\equiv\)}
\begin{LitwebMarkedCode}
\LitwebCodeLine{}{\LitwebCodeDefinition{4}}fn main() \LitwebLeftBrace{}
\LitwebCodeLine{}{}    \(\langle\)\texttt{Compute_value #1}\ \LitwebCodeLocation{1}{2}{2}\(\rangle\)
\LitwebCodeLine{}{\LitwebCodeDefinition{5}}    let path = r"C:\LitwebBackslash{}very\LitwebBackslash{}long\LitwebBackslash{}path\LitwebBackslash{}whose\LitwebBackslash{}contents\LitwebBackslash{}make\LitwebBackslash{}this\LitwebBackslash{}source\LitwebBackslash{}line\LitwebBackslash{}wrap\LitwebBackslash{}without\LitwebBackslash{}becoming\LitwebBackslash{}active\LitwebBackslash{}LaTeX";
\LitwebCodeLine{}{}    println!("\LitwebLeftBrace{}path\LitwebRightBrace{}");
\LitwebCodeLine{}{}\LitwebRightBrace{}
\end{LitwebMarkedCode}
\LitwebSection{litweb-1-2}{2}
\LitwebCodeTitle{\(\langle\)\texttt{Compute\_value\ \#1}\ \hyperlink{litweb-1-2}{2}\(\rangle\) \(\equiv\)}
\begin{LitwebMarkedCode}
\LitwebCodeLine{}{\LitwebCodeDefinition{1}}struct Graph \LitwebLeftBrace{}
\LitwebCodeLine{}{\LitwebCodeDefinition{7}}    value_1: i32,
\LitwebCodeLine{}{}\LitwebRightBrace{}
\LitwebCodeLine{}{}
\LitwebCodeLine{}{\LitwebCodeDefinition{6}\LitwebCodeDefinition{8}\LitwebCodeDefinition{9}}fn squared_distance(x: i32, y: i32) -> i32 \LitwebLeftBrace{}
\LitwebCodeLine{\LitwebCodeUse{8}\LitwebCodeUse{8}\LitwebCodeUse{9}\LitwebCodeUse{9}}{}    x * x + y * y
\LitwebCodeLine{}{}\LitwebRightBrace{}
\end{LitwebMarkedCode}
\LitwebSeeAlso{See also section \hyperlink{litweb-1-3}{3}.}
\LitwebSeeAlso{This code is used in section \hyperlink{litweb-1-1}{1}.}
\LitwebSection{litweb-1-3}{3}
\LitwebCodeTitle{\(\langle\)\texttt{Compute\_value\ \#1}\ \hyperlink{litweb-1-2}{2}\(\rangle\) +\(\equiv\)}
\begin{LitwebMarkedCode}
\LitwebCodeLine{\LitwebCodeUse{1}}{\LitwebCodeDefinition{3}\LitwebCodeDefinition{2}}fn graph_value(graph: &Graph) -> i32 \LitwebLeftBrace{}
\LitwebCodeLine{\LitwebCodeUse{2}\LitwebCodeUse{7}}{}    graph.value_1
\LitwebCodeLine{}{}\LitwebRightBrace{}
\end{LitwebMarkedCode}
\LitwebSeeAlso{This code is used in section \hyperlink{litweb-1-1}{1}.}
\LitwebMiniIndexEnd
\section*{Identifier Index}
\begin{description}
\item[\texttt{Graph}] \hyperlink{litweb-1-1}{1}, \LitwebIdentifierDefinition{\hyperlink{litweb-1-2}{2}}, \hyperlink{litweb-1-3}{3}
\item[\texttt{graph}] \LitwebIdentifierDefinition{\hyperlink{litweb-1-3}{3}}
\item[\texttt{graph\_value}] \LitwebIdentifierDefinition{\hyperlink{litweb-1-3}{3}}
\item[\texttt{main}] \LitwebIdentifierDefinition{\hyperlink{litweb-1-1}{1}}
\item[\texttt{path}] \LitwebIdentifierDefinition{\hyperlink{litweb-1-1}{1}}
\item[\texttt{println}] \hyperlink{litweb-1-1}{1}
\item[\texttt{squared\_distance}] \LitwebIdentifierDefinition{\hyperlink{litweb-1-2}{2}}
\item[\texttt{value\_1}] \hyperlink{litweb-1-1}{1}, \LitwebIdentifierDefinition{\hyperlink{litweb-1-2}{2}}, \hyperlink{litweb-1-3}{3}
\item[\texttt{value\_2}] \hyperlink{litweb-1-1}{1}
\item[\texttt{x}] \LitwebIdentifierDefinition{\hyperlink{litweb-1-2}{2}}
\item[\texttt{y}] \LitwebIdentifierDefinition{\hyperlink{litweb-1-2}{2}}
\end{description}
\end{document}
